\section{\textbf{Limitations}}

FactorJEPA demonstrates that structured predictive channels can improve world modeling in dense, heterogeneous, and partially observed urban scenes. The present evidence nevertheless leaves several important questions open, particularly around interaction grounding, semantic factorization, long-horizon evaluation, geographic transfer, and full-scale training.

\paragraph{\textbf{\emph{(i) Interaction grounding remains the principal open challenge.}}} Our interaction targets combine frozen DINOv2 regions, temporal association, relative motion, visibility, proximity, and reliability weighting. This design provides scalable supervision across 1,000 hours of video and captures a tractable subset of the interaction cues that matter for prediction. However, urban interactions are not directly segmentable visual objects: they are temporally extended relations shaped by anticipation, yielding, hesitation, shared right of way, informal signaling, and partially observed actors. Nearby agents may move independently, whereas distant or occluded agents may still influence the future. The current pathway therefore emphasizes sparse pairwise coupling and does not yet fully represent \textbf{\emph{(i) group-level motion}}, \textbf{\emph{(ii) multi-agent negotiation}}, or \textbf{\emph{(iii) temporally persistent interaction events}}. The independently annotated audit split ensures that the training-time interaction generator is not reused as the headline evaluator, while reliability weighting reduces the influence of uncertain targets. Even so, richer event-level annotation, trajectory-aware grounding, and higher-order relational representations could provide more complete supervision. Accordingly, FactorJEPA should be interpreted as learning \textbf{\emph{interaction-sensitive predictive channels}}, while richer semantic, intentional, and causal interaction modeling remains open.

\paragraph{\textbf{\emph{(ii) Factor semantics are anchored at the channel level.}}} Layout, agent, visibility, and interaction targets are produced by a fixed DINOv2-based pipeline rather than exhaustive human annotation. This enables supervision at dataset scale, but the resulting targets may still reflect missed small agents, fragmented tracks, uncertain boundaries, or reduced reliability under severe occlusion, blur, poor illumination, and camera motion. Some distinctions are also inherently contextual: parked vehicles, temporary barriers, vendors, crowds, and encroachments may behave as persistent layout in one scene and dynamic entities in another. Factor-specific pathways, heads, and separation losses provide \textbf{\emph{semantic anchoring and functional channel separation}}, but do not imply coordinate-level identifiability, complete statistical independence, or a unique decomposition of the future latent. Alternative teachers, taxonomies, or factor ranks may yield different yet comparably predictive partitions. Multi-teacher agreement, teacher-free factor discovery, and intervention-based equivalence tests are therefore natural extensions.

\paragraph{\textbf{\emph{(iii) The evaluation targets predictive structure rather than complete causal understanding.}}} The four primary diagnostics have deliberately bounded interpretations. Future-frame L1 measures fidelity in the frozen target-encoder space; Mask-ratio slope measures robustness under controlled evidence removal; Motion cosine measures linearly accessible motion rather than all motion information encoded by the model; and Causal L1 evaluates consistency between predicted and target-encoder changes under independently specified interventions. Causal L1 therefore measures \textbf{\emph{intervention sensitivity}}, not causal identification or unrestricted counterfactual reasoning. The experiments also emphasize short-horizon prediction from prerecorded video. Longer rollouts may accumulate identity, layout, and interaction errors, while action-conditioned forecasting, closed-loop planning, active perception, and online adaptation remain untested. The current results thus establish FactorJEPA as a \textbf{\emph{structured predictive representation}}; extending it to long-horizon and closed-loop world modeling remains future work.

\paragraph{\textbf{\emph{(iv) Geographic breadth does not imply universal coverage.}}} DENSEWORLD spans 22 Indian cities, three capture modes, and substantial variation in density, infrastructure, illumination, weather, road structure, and traffic composition. It should nevertheless be viewed as \textbf{\emph{one large-scale realization}} of populous, crowded, and chaotic Global South urban environments rather than an exhaustive characterization of the Global South. Mobility systems, infrastructure, road conventions, climate, and social negotiation differ across South Asia, Africa, Latin America, Southeast Asia, and the Middle East. The corpus also emphasizes outdoor urban public spaces; rural roads, indoor crowds, industrial environments, disaster settings, and other interaction regimes remain outside its present scope. Cross-region transfer, geographically held-out evaluation, and broader collection are needed to test how far the learned factorization generalizes.

\paragraph{\textbf{\emph{(v) Full-scale training and visual decoding remain incomplete.}}} Full 115k-clip training is reported for the 1B backbone, whereas the 2B model is evaluated under the matched stratified 10k regime. The strong cross-scale rank correlations indicate stable relative method behavior, but they do not replace a full-data 2B experiment. It therefore remains open whether larger-scale training further strengthens the predictive channels, changes their depth-wise allocation, or modifies the observed prediction--motion trade-off. The Cosmos-initialized latent-to-RGB decoder also introduces a separate rendering stage. Oracle controls isolate forecast-induced error from decoder reconstruction error, but rendering may still exhibit blur, deterministic averaging, missed fine-grained agents, or uncertainty collapsed into a single future. Visual plausibility should therefore be interpreted alongside latent-space and oracle-controlled evaluation rather than as standalone evidence of predictive correctness.

These limitations motivate three immediate directions: \textbf{\emph{(i) richer interaction grounding and event-level supervision}}, \textbf{\emph{(ii) higher-order and temporally persistent relational modeling}}, and \textbf{\emph{(iii) broader geographic, long-horizon, and full-scale validation}}.